\maketitle
\head{本版编排与收录说明}
\noindent\textbf{这是一份按内容顺序整理的中文工作译稿，不是全部已上传材料的完结全译本。}第I、II章及第III章的两个已出版分册正文已经连续收入；第0章当前收入§§1--18及§19的19.0--19.4；第IV章主体、第0章19.5--19.10及§§20--23尚未连续译入；第V卷英文预备笔记只收入已译出的部分。未译内容没有用原文扫描页、提要或空白标题代替译文。

\section*{正文次序}
本版不再按上传时间、译稿批次或独立PDF页序排列，而按原著内容编号组织：导言；第0章预备知识；第I章概形的语言；第II章若干类态射的初等整体研究；第III章凝聚层的上同调研究；后续待译范围说明；独立编排的第V卷预备笔记已译部分；最后为各原分册的书末资料。第III章第一部分的§§1--5与第二部分的§§6--7直接衔接，不再被第一部分的索引隔开。

第0章原来分载于第一卷、第三卷第一部分、第四卷第一部分。本版将已译的§§1--7、§§8--13、§§14--18和§19已译部分归并在一起，但保留 $0_{\mathrm I}$、$0_{\mathrm{III}}$、$0_{\mathrm{IV}}$ 等出处记号。它们标示原出版分册，并非重复的三章。这样的安排依内容编号，而不声称每一节的全部前提均已在前文出现；原著跨章引用仍按原号保留。

\section*{实际收录范围}
\begin{longtable}{@{}p{.25\linewidth}p{.67\linewidth}@{}}
\toprule 材料 & 已收入的连续译文及状态\\\midrule\endhead
第一卷（1960） & 导言；第0章§§1--7；第I章§§1--10，止于（10.15.7）；原书书末资料。承接既有V2校订源码。\\
第二卷（1961） & 第II章§§1--8，止于（8.14.14）；原书参考文献、索引、目录、勘误清单一。其书末资料沿用上一合订版依据底本恢复的译排来源。\\
第三卷第一部分（1961） & 第0章§§8--13；第III章§§1--5，止于（5.5.2）；原书书末资料。\\
第三卷第二部分（1963） & 第III章§§6--7，止于（7.9.14），连同其证明、记号索引、术语索引、原目录及勘误清单二。第三版已补齐（6.7.1）至本分册末尾；本版保留。\\
第四卷第一部分（1964） & 第0章补篇章首说明、§§14--18及§19的19.0--19.4，止于（19.4.7）及其证明。本版新增（19.0.1）--（19.4.7）；19.5以后尚未译入。\\
第V卷英文预备笔记 & \texttt{EGAV1-2.ps} 的六页可读材料；\texttt{EGAVaiello.pdf} 的卷首说明及“超平面截面与锥投影”§§1--6。余下内容不计入已译范围。\\
\bottomrule
\end{longtable}

\section*{译文、勘误与核查的关系}
译文保留原陈述与证明的详细程度。原著省略证明之处仍明确保留省略；原稿未写完、英译编者标记无法辨认或待查之处，不凭通常知识补造“原文”。原注、英译编者注与本次译校注使用不同前缀。本版新增译校注以F6开头；第五版E5、第四版D4、第三版C3及更早批次编号保留。对可确定的排印调整，在相应注释中列出原式及处理依据；对于证明缺步，区分“该推理不能成立”和“命题结论已被反例推翻”。

此前收入的第III卷第二部分书末《勘误与增补（清单二）》是\textbf{作者已经发表的勘误}，不应当作本次的新发现。其中涉及第一卷进制完备性、对角态射，以及第二卷若干证明的修正，已在对应正文旁添加可跳转的提示。原正文与早期译注仍保留，避免把历史原文与修订后的表述静默混合。

第V卷上传材料的标题明确写为预备笔记的翻译与编辑，英译编者为 Piotr Blass 与 Joseph Blass；它们包含重组建议、未完成句子及待补证明，不按正式刊行的第I--IV章同等状态标示。与合并英文稿重合的PS载体不重复计作新译文。

\section*{第六版的新增范围与修订方式}
本版连续译入第0章§19的引言、形式满态射与形式单态射、形式投射模、形式光滑代数、形式光滑性的初步判据，即（19.0.1）至（19.4.7），保留其间子编号引理、证明与原注。对应第四卷第一部分分册印刷第69--89页（上传PDF第68--88页）。下一接续点为19.5小节及（19.5.1）；本版没有将已读取的19.5开头算作译完。

作者已在第四卷第四部分的正式勘误中要求删除（19.3.12）。本版为保存初版对照而译出该条，但在陈述前后显著标出撤销状态，并将正式勘误选译到书末；它不作为未经修订的有效命题供引用。核查报告还区分（19.3.10.1）按原列假设确有反例、（19.3.5）有限乘积结论成立而原证明错误，以及（19.4.1）可以补全的选取步骤。不同类别不混作“错误定理”。

此前第0章§§1--18、第I--III章及预备笔记正文保留，不因本轮续译而改写，也没有宣称本次已重新逐句审定旧稿。只有新增范围及本轮直接核对的记号、正式勘误和证明步骤纳入第六版的新核查。

\section*{版式与使用}
采用 \texttt{\string\documentclass[UTF8,a4paper]\{ctexbook\}}，由文档类决定默认中文字体、字号、版心、标题与双面分页。不另设页边距，不用自定义页眉。为保留原著结构，仅保留原章号、节号与公式号；默认双面排版所需的空白偶数页不删除。前置部分使用罗马数字页码，其后连续编号；总目录指向本合订本页面，各原著目录的页码仍指原版。

PDF可用于对照研读；可独立编译的单文件TeX与分文件工程同时提供。旧稿内容的编排归并与技术校验，不等于本次将旧稿的每一句译文或每一个证明重新独立审定。累计核查报告单独交付，历史核查意见带有批次说明；对同一事项，本版直接核得的作者勘误和明确论证须与历史意见对照阅读，不能将历史“待核”标签当作最新状态。
\tableofcontents
